\section{Final remarks}\label{sec:conclusion}

This paper is centered on a finite-state theorem: for first-entry
residue events of a Markov hole, the depth-\(m\) Bayes error is exactly
a killed-matrix path sum with residue-dependent terminal ambiguity
weights.  From that identity one recovers the escape rate
\(\gamma_H(\varphi)\), the residue-dependent quasistationary ambiguity
amplitude \(\xi_H^{\mathsf T}b_r\), and automatic positivity of this
amplitude for every residue.  The latter is forced by mixing: a survivor
loop of any sufficiently large length may be concatenated with a fixed
exit path, so every recurrent survivor state realizes first entries in
every residue class.  The rest of the paper is intentionally secondary
to that theorem.  The pressure discussion records a standard survivor
corollary needed for iid collision estimates, while the cyclotomic
resolvent discussion is kept at the level it proves: an algebraic
repackaging of the same killed path sum as a finite-matrix generating
function whose possible scalar poles lie in the reciprocal cyclotomic
lift of the killed spectrum, with visibility determined by spectral
non-cancellation.  In the one-step Markov setting these objects are
completely explicit, and the doubling-hole example yields the closed-form
constant \(2/5\).

The standard survivor R\'enyi pressure identity
\(s\mapsto sP_{Y_H}(\varphi)-P_{Y_H}(s\varphi)\) packages the
iid collision behaviour of surviving samples into open-system thermodynamic
formalism.  At \(s=2\) it determines the birthday scale of the
conditioned survivor law under independent sampling, with the ambient
two-sample exponent splitting
as \(\chi_H = 2\gamma_H(\varphi)+h_{2,H}\).  If the survivor has at
least two states, the strict spectral inequality
\[
  h_{t,H}>\frac{t}{s}h_{s,H}\qquad(1\leq s<t)
\]
implies in particular \(h_{3,H}>\frac32h_{2,H}\).  It therefore makes
the Chen--Stein third-moment term exponentially negligible at the
birthday scale, with no additional overlap hypothesis.  The one-state
survivor has \(h_{s,H}=0\) for every \(s\) and hence has no exponential
birthday regime.

Two directions remain natural.  One is an operator-theoretic extension
of the finite-observation theorem with the functional-analytic
hypotheses proved directly rather than by finite reduction.  The other
is to understand whether finer residue observables or multidimensional
holes make larger pieces of the finite-state open spectrum visible than
the scalar reciprocal cyclotomic resolvents considered here.
